\documentclass[10pt,a4paper]{amsart}
\usepackage[margin=19mm]{geometry}
\usepackage{amsmath,amssymb,mathtools}
\usepackage[scr=rsfs,cal=euler]{mathalfa}
\usepackage{xfrac}
\usepackage[unicode,hidelinks,bookmarks=false]{hyperref}
\newcommand{\ie}{i.e.\ }
\newcommand{\cf}{cf.\ }
\newcommand{\resp}{respectively\ }
\newcommand{\Subsection}{\subsection}
\providecommand{\nobreakdash}{\nobreak\hbox{-}}
\setlength{\emergencystretch}{2em}
\multlinegap0pt
\usepackage{amssymb}
\usepackage{enumerate}
\usepackage{algorithm}\usepackage{algpseudocode}
\usepackage{dsfont}
\usepackage[shortlabels]{enumitem}
\usepackage{subfig}
\newtheorem{lemma}{Lemma}
\newtheorem{proposition}{Proposition}
\usepackage{cleveref}
\crefname{lemma}{Lemma}{Lemmas}
\crefname{proposition}{Proposition}{Propositions}
\newcommand{\bR}{\mathbb{R}}
\newcommand{\cB}{\mathcal{B}}
\newcommand{\cC}{\mathcal{C}}
\DeclareMathOperator{\dist}{dist}
\DeclareMathOperator{\dom}{dom}
\DeclareMathOperator{\infconv}{\mathbin{\square}}
\title{All roads lead to Rome: \\ Path-following Augmented Lagrangian Methods \\ via Bregman Proximal Regularization}
\author{Emanuel Laude}
\date{Source: arXiv:2602.15710v1 (CC BY 4.0)}
\begin{document}
\maketitle

\noindent\emph{Source and licence.} All roads lead to Rome: \\ Path-following Augmented Lagrangian Methods \\ via Bregman Proximal Regularization. Version arXiv:2602.15710v1 (CC BY 4.0), distributed by arXiv under the Creative Commons Attribution 4.0 International licence, http://creativecommons.org/licenses/by/4.0/, as recorded in the arXiv licence metadata for this exact version and shown on the abstract page https://arxiv.org/abs/2602.15710v1. Full licence notice: \texttt{licenses/arxiv-2602.15710-CC-BY-4.0.txt}.\par\smallskip

\noindent\textit{Editorial scope.} Counted unit: the article's proposition proposition (source0.tex lines 1238--1245). The article's complete original proof of it is delivered with it (source0.tex lines 1246--1293, one \texttt{proof} environment of 48 lines). No mathematics is written, repaired or re-derived: the statement and the proof are the article's own lines, in the article's own notation, and no retained line is substituted or re-flowed.  Retained uncounted as the source-local context the proof needs: lines 1166--1171.  Nothing was omitted from any retained span, so the package carries no omission record.  No retained line contains a citation, so the wrapper carries no bibliography and no bibliography entry.  No retained line contains a figure environment, an image-inclusion command or drawing code, so no image is delivered and the package declares no asset dependency.  Editorial formatting only, in addition: an AMS \texttt{amsart} preamble that restates the article's own declarations and macros used by the retained lines, namely the environments \texttt{lemma}, \texttt{proposition} on separate counters, together with the standard packages those lines need.  The retained environments print in this document as Lemma 1, Proposition 1 in order of appearance, whereas the article prints the same items with the numbering of its own full document.  The complete native source of the article as distributed in the arXiv e-print for this exact version is bundled unchanged as \texttt{sources/194816/source0.tex}.
\begin{lemma} \label{thm:ergodic_pd_fixed}
    Let $\psi \in \Gamma_0(\bR^n)$ and $\phi \in \Gamma_0(\bR^m)$ be Legendre. Then it holds that for any $x \in \dom \psi$ and $y \in \dom \phi$
\begin{align}
    L(s^{k}, y) -L(x, y^{k+1})&\leq \tfrac{1}{\sigma_k}\big(D_\Phi(z,z^k) -  D_\Phi(z,z^{k+1}) - (1-\rho_k) D_\Phi(p^k,z^{k})\big).
\end{align}
\end{lemma}

\begin{proposition}[asymptotic feasibility for conic constraints]
Let $\psi \in \Gamma_0(\bR^n)$ and $\phi \in \Gamma_0(\bR^m)$ be Legendre. Let $g=\delta_{-\cC}$ for a closed and convex cone $\mathcal{C} \subseteq \bR^m$ such that $\cC^* \subseteq \dom \phi$. Let $(x^\star, y^\star) \in \dom \Phi$ be a saddle-point of $L$.
Then we have that
$$
\max\{|f(\breve s^K) - f(x^\star)|, \dist(\mathcal{A}(\breve s^K), -\cC)\} \leq \frac{1}{\sum_{k=0}^{K} \sigma_k} \big(D_\psi(x^\star,x^0) + \max_{y \in \cC^* \cap B_R} D_\phi(y, y^0)\big),
$$
where $R:=2\|y^\star\| + 1$.
\end{proposition}

\begin{proof}
Let $y \in \cC^*$. Since $\cC^*$ is convex and $\{y^{k}\}_{k=0}^\infty \subset \cC^*$ we have that $\breve y^{K} \in \cC^*$ and since $\mathcal{A}(x^\star) \in -\cC$ we obtain that $\langle \breve y^{K}, \mathcal{A}(x^\star) \rangle \leq 0$. Thus we have
\begin{align*}
L(\breve s^K, y) - L(x^\star, \breve y^K) &= f(\breve s^K) - f(x^\star) + \langle y, \mathcal{A}(\breve s^K)\rangle - \langle \breve y^{K}, \mathcal{A}(x^\star) \rangle \\
&\geq f(\breve s^K) - f(x^\star) + \langle y, \mathcal{A}(\breve s^K)\rangle.
\end{align*}
Invoking \cref{thm:ergodic_pd_fixed} we obtain
\begin{align}
\frac{1}{\sum_{k=0}^{K} \sigma_k} \big(D_\psi(x^0,x^\star) + D_\phi(y, y^0)\big) &\geq L(\breve s^K, y) - L(x^\star, \breve y^K) \notag \\
&\geq f(\breve s^K) - f(x^\star) + \langle y, \mathcal{A}(\breve s^K)\rangle\label{eq:bound_primal_dual}.
\end{align}
For $y =0$ we have that
\begin{align} \label{eq:ergodic_simple}
    f(\breve s^K) - f(x^\star) \leq \frac{1}{\sum_{k=0}^{K} \sigma_k} \big(D_\psi(x^0,x^\star) + D_\phi(0, y^0)\big).
\end{align}
Let $R:= 2\|y^\star\| + 1$.
Taking the maximum over $y \in \cC^* \cap B_R$ on both sides of the inequality \cref{eq:bound_primal_dual} noting that
\begin{align*}
\max_{y \in \cC^* \cap B_R} \langle y, u\rangle = (\delta_{\cC^*} + \delta_{B_R})^*(u)  =(\delta_{-\cC} \infconv (R\|\cdot\|))(u) =R \dist(u,-\cC)
\end{align*}
we obtain
\begin{align}
    \frac{1}{\sum_{k=0}^{K} \sigma_k} \big(D_\psi(x^0,x^\star) + \max_{y \in \cC^* \cap B_R} D_\phi(y, y^0)\big) \geq f(\breve s^K) - f(x^\star) + R \dist(\mathcal{A}(\breve s^K), -\cC).\label{eq:ergodic_1}
\end{align}
Since $(x^\star,y^\star)$ is a saddle-point of $L$ we have for all $x \in \bR^n$ and $y \in \cC^*$ that
$$
f(x) + \langle \mathcal{A}(x), y^\star \rangle-f(x^\star) - \langle \mathcal{A}(x^\star), y\rangle= L(x, y^\star) - L(x^\star, y) \geq 0
$$
and thus for $y=0$ and $x=\breve s^K$ since $y^\star \in \cC^*$
\begin{align} 
0 &\geq f(x^\star)-f(\breve s^K) -\langle y^\star, \mathcal{A}(\breve s^K) \geq f(x^\star)-f(\breve s^K) - \max_{y \in \cC^*\cap \|y^\star\|\cB} \langle y, \mathcal{A}(\breve s^K)\rangle \notag \\
&= f(x^\star)-f(\breve s^K) - \|y^\star\|\dist(\mathcal{A}(\breve s^K),-\cC) \label{eq:ergodic_2}
\end{align}
Summing \cref{eq:ergodic_1} and \cref{eq:ergodic_2} since $R=2\|y^\star\|+1$ we obtain
\begin{align}
\dist(\mathcal{A}(\breve s^K),-\cC) &\leq (\|y^\star\| + 1)\dist(\mathcal{A}(\breve s^K),-\cC) \notag \\
&\leq \frac{1}{\sum_{k=0}^{K} \sigma_k} \big(D_\psi(x^0,x^\star) + \max_{y \in \cC^* \cap B_R} D_\phi(y, y^0)\big) \label{eq:final_primal_dual}.
\end{align}
From \cref{eq:ergodic_2} and \cref{eq:final_primal_dual} we obtain
\begin{align*}
    f(x^\star)-f(\breve s^K) &\leq \|y^\star\| \dist(\mathcal{A}(\breve s^K),-\cC) \\
    &\leq \frac{1}{\sum_{k=0}^{K} \sigma_k} \big(D_\psi(x^0,x^\star) + \max_{y \in \cC^* \cap B_R} D_\phi(y, y^0)\big).
\end{align*}
Together with \cref{eq:ergodic_simple} we obtain
\[
|f(x^\star)-f(\breve s^K)| \leq \frac{1}{\sum_{k=0}^{K} \sigma_k} \big(D_\psi(x^0,x^\star) + \max_{y \in \cC^* \cap B_R} D_\phi(y, y^0)\big) \qedhere
\]
\end{proof}

\end{document}
